\chapter{TRNet Architecture and Training Ablations}
\label{app:trnet_modifications}

This appendix reports the ablation experiments that informed the encoder, decoder and
loss-function choices used in all main-chapter TR-MISR experiments
(Section~\ref{ssec:experimental_setup}).
Three design choices were evaluated independently against the chosen baseline:
First, feeding the encoder a single fluorescence channel versus concatenating an averaged
reference image as a second channel (Section~\ref{sec:adaptation});
secondly, replacing the PixelShuffle decoder with a transposed-convolution (deconvolution) decoder;
and thirdly, replacing the combined L1+Fourier loss with the Fourier-only loss.
All other hyperparameters are held fixed at the values described in
Section~\ref{ssec:experimental_setup}.

%% -------------------------------------------------------
\section{Encoder: Single-channel vs.\ Added Averaged Reference Image}
\label{app:encoder_ablation}

The original TR-MISR encoder takes a two-channel input, concatenating each
low-resolution frame with a reference image that acts as an implicit
co-registration anchor~\cite{TR-MISR-original}. As argued in
Section~\ref{sec:adaptation}, in SOFI the frame stack is co-registered by
construction, so this reference is expected to be redundant. This ablation tests
that claim empirically by comparing two otherwise identical 1-head, 1-depth, 
TR-MISR models evaluated at a 20-frame budget:
the chosen single-channel model, and a two-channel variant whose second channel
is the per-stack averaged reference image. We chose average over median as the background 
noise  will get averaged out and only the diffraction limited structure remains.

The comparison is paired: both models are evaluated on the identical 480-sample
test set and scores are matched by condition and sample identity. Following the
procedure in Appendix~\ref{app:significance}, we report the mean
resolution-scaled Pearson (RSP) difference $\overline{d}$, the two-sided
Wilcoxon signed-rank $p$-value, and the matched-pairs rank-biserial correlation.

% \begin{table}[htbp]
%   \centering
%   \caption{RSP per SNR condition for the single-channel and averaged-reference
%            encoders at a 20-frame budget. $\overline{d}$ is the mean RSP
%            difference (single\,$-$\,avg.\ ref.); all rows are significant under
%            the paired Wilcoxon test (per-condition $p = 1.8\times10^{-12}$,
%            rank-biserial $=-1.00$, i.e.\ the single-channel model wins on every
%            one of the 40 samples). Bold marks the better value per row.}
%   \label{tab:ablation_encoder}
%   \begin{tabular}{lccc}
%     \toprule
%     \textbf{Condition} & \textbf{Single channel} & \textbf{Avg.\ reference} & $\overline{d}$ \\
%     \midrule
%     01 fast p100   & \textbf{0.9048} & 0.8978 & $+0.0070$ \\
%     02 fast p80    & \textbf{0.9048} & 0.8979 & $+0.0069$ \\
%     03 fast p60    & \textbf{0.9100} & 0.9016 & $+0.0084$ \\
%     04 fast p40    & \textbf{0.9079} & 0.8968 & $+0.0111$ \\
%     \midrule
%     05 medium p100 & \textbf{0.9033} & 0.8867 & $+0.0166$ \\
%     06 medium p80  & \textbf{0.9066} & 0.8887 & $+0.0178$ \\
%     07 medium p60  & \textbf{0.9056} & 0.8838 & $+0.0218$ \\
%     08 medium p40  & \textbf{0.9103} & 0.8839 & $+0.0264$ \\
%     \midrule
%     09 slow p100   & \textbf{0.8952} & 0.8558 & $+0.0394$ \\
%     10 slow p80    & \textbf{0.8939} & 0.8497 & $+0.0442$ \\
%     11 slow p60    & \textbf{0.9000} & 0.8494 & $+0.0506$ \\
%     12 slow p40    & \textbf{0.8871} & 0.8188 & $+0.0684$ \\
%     \midrule
%     \textbf{Aggregate} & \textbf{0.9025} & 0.8759 & $+0.0266$ \\
%     \bottomrule
%   \end{tabular}
% \end{table}

The single-channel encoder is significantly better in every condition and on
every paired sample. Pooled over all 480 samples it gains $\overline{d} = +0.0266$
RSP ($p \approx 2.3\times10^{-80}$, rank-biserial $=-1.00$). The margin grows monotonically
as the imaging conditions worsen, from $+0.007$ RSP at the easiest fast/high-photon
setting to $+0.068$ RSP at the hardest 12\_slow\_p40 condition. Concatenating an
averaged reference image therefore does not help and actively hurts: the reference
adds a redundant, low-fluctuation channel that dilutes the blinking statistics the
fusion relies on, with the largest cost precisely where the fluctuation signal is
weakest. This confirms the design choice of a single-channel encoder
(Section~\ref{sec:adaptation}).

%% -------------------------------------------------------
\section{Decoder: PixelShuffle vs.\ Transposed Convolution}
\label{app:decoder_ablation}

The SOFI super-resolution task requires upscaling LR from $N \times N$ to HR result of size 
$2*N-7 \times 2*N-7$.
Two standard decoder approaches were compared:

\begin{itemize}
  \item \textbf{PixelShuffle} (sub-pixel convolution~\cite{shi2016pixelshuffle}):
   Originaly used in the TR-MISR paper~cite\cite{TR-MISR-original}. It learns
    to rearrange channel dimensions into spatial resolution. The non-integer scale is handled
    by upsampling to an intermediate size followed by a centre crop to match the $2N-7$.
  \item \textbf{Transposed convolution} (\texttt{ConvTranspose2d}): Is used in MISRGRU models from 
  RESURF~\cite{resurf2024} learned upsampling via
    strided deconvolution. Known to produce checkerboard artefacts when the kernel size is
    not divisible by the stride~\cite{odena2016deconvolution}.
\end{itemize}

%% -------------------------------------------------------
\section{Loss Function: L1+Fourier vs.\ Fourier Only}
\label{app:loss_ablation}

The Fourier loss~\cite{resurf2024} penalises errors in the amplitude and phase spectra of
the prediction, encouraging recovery of high-frequency detail that an L1 loss would
suppress via mean regression. A combined L1+Fourier loss retains the spatial stability of
L1 while adding spectral supervision. The ablation tests whether the L1 term is redundant
or provides complementary signal.

%% -------------------------------------------------------
\section{Results}
\label{app:ablation_results}

Table~\ref{tab:ablation_aggregate} reports aggregate metrics over all 480 test samples
(12 SNR conditions $\times$ 40 samples each). Table~\ref{tab:ablation_per_snr} gives
the RSP breakdown per SNR condition.

\begin{table}[htbp]
  \centering
  \caption{Aggregate metrics for the three configurations over 480 test samples.
           Bold marks the best value per column.}
  \label{tab:ablation_aggregate}
  \begin{tabular}{lcccc}
    \toprule
    \textbf{Model} & \textbf{RSP}  \\
    \midrule
    L1+Fourier, PixelShuffle (baseline) & \textbf{0.8789} \\
    Fourier only, PixelShuffle          & 0.8757          \\
    L1+Fourier, Deconv                  & 0.8770          \\
    \bottomrule
  \end{tabular}
\end{table}

\begin{table}[htbp]
  \centering
  \caption{RSP per SNR condition. Conditions are ordered by blinking speed
           (fast/medium/slow) and photon level (100\%--40\%). Bold marks the best
           value per row.}
  \label{tab:ablation_per_snr}
  \begin{tabular}{lccc}
    \toprule
    \textbf{Condition} & \textbf{L1+F, PS (baseline)} & \textbf{F only, PS} & \textbf{L1+F, Deconv} \\
    \midrule
    01 fast p100  & \textbf{0.9013} & 0.8987 & 0.8995 \\
    02 fast p80   & \textbf{0.8996} & 0.8984 & 0.8988 \\
    03 fast p60   & 0.9022          & \textbf{0.9024} & 0.9021 \\
    04 fast p40   & 0.8950          & \textbf{0.8998} & 0.8964 \\
    \midrule
    05 medium p100 & \textbf{0.8885} & 0.8881 & 0.8872 \\
    06 medium p80  & 0.8900          & \textbf{0.8915} & 0.8892 \\
    07 medium p60  & 0.8845          & \textbf{0.8872} & 0.8839 \\
    08 medium p40  & 0.8843          & \textbf{0.8861} & 0.8842 \\
    \midrule
    09 slow p100   & \textbf{0.8579} & 0.8548 & 0.8545 \\
    10 slow p80    & \textbf{0.8536} & 0.8483 & 0.8498 \\
    11 slow p60    & \textbf{0.8556} & 0.8453 & 0.8510 \\
    12 slow p40    & \textbf{0.8345} & 0.8073 & 0.8269 \\
    \midrule
    \textbf{Aggregate} & 
     & 0.8757 & 0.8770 \\
    \bottomrule
  \end{tabular}
\end{table}

\subsection{Decoder comparison}

The Deconv decoder performs comparably to PixelShuffle at high-SNR conditions (fast
blinking, high photon counts) but falls behind consistently at medium and slow blinking.
At the most challenging condition (12\_slow\_p40), the PixelShuffle baseline achieves
RSP~$=0.835$ versus $0.827$ for Deconv, a gap of $0.008$ that represents a meaningful
difference in structural fidelity under low-SNR conditions.

% \todo{comparison_results/eval_trnets_fourier_deconv/multi_comparison_results.json statistical analysis}

\subsection{Loss function comparison}

The Fourier-only loss matches or slightly exceeds the baseline at fast-blinking, high-SNR
conditions (conditions 03--04 and 06--08) but degrades at slow-blinking,
low-photon conditions. At condition 12\_slow\_p40, the Fourier-only model achieves
RSP~$=0.807$ versus $0.835$ for the baseline, a gap of $0.027$.
 The gap between baseline and Fourier-only grows with decreasing SNR.

The L1 term therefore provides complementary supervision: it penalises spatial
misregistration directly, stabilising the model under challenging imaging conditions
where Fourier-based supervision alone is insufficient.

